% HD-G32. GPT-6 Astra Pro /root/g34; inherited effort not exposed.
% 21 November 2026. Actual primitive normalizer reduction in degrees15--35.
\section{HD-G32: the actual quasi-simple core in the remaining primitive case}
\label{sec:primitive-quasisimple-core}

Retain the actual natural sixfold source of HD-G31, over
\(K=\mathbb Q(i)\) or
\(K=\mathbb Q(\sqrt{-d})\), \(d=2,7,11,19,43,67,163\).
Its exact coefficient character field is \(K\), its polarized
period rank is nine, and its coefficient group has no nonidentity
element of moving rank one or two. Let \(N=N_3\) be the normal
subgroup of the original \(G\subset GL_m(K)\) generated by all
its actual moving-rank-three elements. In the absolutely irreducible,
complex-primitive branch, HD-G31 and HD-R14 give
\(15\le m\le35\) and exclude every nontrivial tensor decomposition,
including permutations of tensor factors.

In this remaining branch there is an actual characteristic
quasi-simple subgroup \(S\triangleleft N\), absolutely irreducible
on the original \(K^m\). In particular the original \(G\) normalizes
\(S\). Writing \(Q=G\cap K^*I\), actual conjugation gives an injection
\[
      G/Q\longrightarrow\operatorname{Aut}(S/Z(S))
      \quad\hbox{whose image contains }\operatorname{Inn}(S/Z(S)).
 \tag{G32.1}
\]
The representation of \(S\) is faithful, its character field is
contained in \(K\), and
\[
                  |Z(S)|\mid\gcd(m,|\mu(K)|).
 \tag{G32.2}
\]
We prove these assertions, including the actual matrix and corner
maps, before stating the precise remaining finite check. No
exclusion of all the remaining quasi-simple cases is asserted.

\paragraph{1. Actual normal restrictions and a minimal projective subgroup.}
Put \(Z=Z(N)\). Schur's lemma and the original \(K\)-matrices give
\(Z=N\cap K^*I\subset\mu(K)I\). Every nonscalar normal subgroup
\(H\triangleleft N\) is absolutely irreducible on the same space.
Indeed finite averaging decomposes the restriction into isotypic
summands, permuted transitively by \(N\). More than one summand
would be a nontrivial system of imprimitivity. Thus there is one
type \(W\), and the literal evaluation map is
\[
 W\otimes_{\mathbb C}\operatorname{Hom}_H(W,V_{\mathbb C})
       \longrightarrow V_{\mathbb C},\qquad w\otimes f\longmapsto f(w).
 \tag{G32.3}
\]
It is an isomorphism by the isotypic decomposition. The actual
image algebra \(\mathbb C[H]\) is
\(\operatorname{End}(W)\otimes I\). For an element normalizing
this algebra, its transformed matrix units give an inner
implementation on \(W\); removing that implementation commutes
with all matrix units and therefore acts on the multiplicity
factor. Thus (G32.3) is an actual fixed tensor decomposition for
\(N\). Its multiplicity must be one by HD-R14, unless \(\dim W=1\),
which would make \(H\) scalar. This proves the assertion.

Choose a minimal nontrivial normal subgroup \(M\) of \(N/Z\),
and let \(H\) be its full inverse image. The elementary finite-group
fact used here is that a minimal normal subgroup is a direct product
of isomorphic simple groups. To recall the argument, it has no
proper nontrivial characteristic subgroup. The automorphic images
of a minimal normal subgroup of it are minimal normal subgroups,
and distinct ones commute and intersect trivially. Their product
is characteristic. A nonabelian minimal normal factor has trivial
centre, because its centre is a proper characteristic subgroup of
that factor and hence normal in the ambient group. Thus in the
nonabelian case these factors form a direct product;
normality in the resulting product then shows that each factor is
simple. In the abelian case it is an elementary abelian group of
one prime exponent, by its characteristic primary and power
subgroups. Transitivity on the factors makes their types equal.
Consequently
\[
 M\simeq C_p^b\quad\hbox{or}\quad M\simeq T^t
       \quad(T\text{ nonabelian simple}).
 \tag{G32.4}
\]
The subgroup \(H\) is nonscalar and is absolutely irreducible by
(G32.3). Hence its centre is exactly \(Z\): every central element
of \(H\) is scalar by Schur, and a scalar lying in \(N\) lies in \(Z\).

\paragraph{2. The elementary-abelian case and its exact trace basis.}
Suppose first that \(H/Z=C_p^b\). The commutator takes values in
\(Z\), and defines the alternating bilinear pairing
\[
 H/Z\times H/Z\longrightarrow\mu_p(K),
                  \qquad (\bar u,\bar v)\longmapsto[u,v].
 \tag{G32.5}
\]
For example \([u^p,v]=[u,v]^p=1\), since \(u^p\in Z\).
The radical is zero because \(Z(H)=Z\). The group \(H\) is
nonabelian, since an irreducible finite abelian group has dimension
one. Thus the pairing has a nonidentity value and \(\mu_p(K)\ne1\).
In all eight fields under consideration this forces \(p=2\).
Successively choosing a pair with nontrivial pairing and taking its
orthogonal complement gives \(H/Z\simeq\mathbb F_2^{2a}\).

Choose one actual representative \(h_v\in H\) for each coset
\(v\), with \(h_0=1\), and put \(B_v=\rho(h_v)\in M_m(K)\).
For \(v\ne0\), choose \(w\) with commutator \(-1\). Conjugation
by \(B_w\) sends \(B_v\) to \(-B_v\), so \(\operatorname{tr}B_v=0\).
The same argument shows
\[
       \operatorname{tr}(B_u^{-1}B_v)=m\,\delta_{u,v}.
 \tag{G32.6}
\]
The averaging operator
\(X\mapsto |H|^{-1}\sum_{h\in H}\rho(h)X\rho(h)^{-1}\)
is projection onto the scalar endomorphisms, by irreducibility.
Its trace on \(\operatorname{End}(V)\) is therefore one. Computing
that trace from the summands, only the scalar cosets contribute,
and gives
\[
 1=|H|^{-1}\sum_{h\in H}|\operatorname{tr}\rho(h)|^2
      =\frac{|Z|m^2}{|H|},\qquad m^2=|H/Z|=2^{2a}.
 \tag{G32.7}
\]
Here the inverse trace is the complex conjugate trace because
all eigenvalues have finite order. Equations (G32.6)--(G32.7)
show that the actual \(K\)-matrices \(B_v\) form a basis of
\(\operatorname{End}_K(K^m)\), and that \(m=2^a\).

Every actual \(x\in N\) normalizes \(H\) and fixes its scalar
centre. Conjugation on this basis is literally
\[
 xB_vx^{-1}=\lambda_v B_{A v},\qquad
 A\in\operatorname{Aut}_{\mathbb F_2}(H/Z),\quad\lambda_v\in\mu(K).
 \tag{G32.8}
\]
For \(v\) fixed by \(A\), multiplication
\(B_uB_v=c(u,v)B_{u+v}\), with actual central scalar \(c(u,v)\),
shows that \(\lambda\) is a character of
\(F=\ker(A-I)\): \(\lambda_{u+v}=\lambda_u\lambda_v\).
Therefore
\[
 |\operatorname{tr}x|^2
   =\operatorname{tr}(\operatorname{Ad}x\mid\operatorname{End}V)
   =\sum_{v\in F}\lambda_v
   \in\{0,|F|\} =\{0,2^{\dim F}\}.
 \tag{G32.9}
\]
This is a calculation in the original matrix algebra, with no
external scalar extension of the finite group.

The only powers of two in \([15,35]\) are16 and32. By HD-G31,
every actual rank-three generator has three nontrivial eigenvalues
\((-1,-1,-1)\), or, only over \(\mathbb Q(i)\), one of
\((i,i,-i)\) and \((-i,-i,i)\). The resulting squared traces are
\[
\begin{array}{c|cc}
 m&(m-6)^2&(m-3)^2+1\\\hline
 16&100&170\\
 32&676&842.
\end{array}                                                   \tag{G32.10}
\]
None is zero or a power of two. This contradicts (G32.9) for
an actual generator of the nontrivial group \(N\). Thus the
first alternative in (G32.4) is excluded.

\paragraph{3. The perfect lift and the actual tensor factors.}
Now \(H/Z=T^t\) with \(T\) nonabelian simple. Put \(E=H'\).
Perfection of \(H/Z\) gives \(H=EZ\), whence
\(E=H'=[EZ,EZ]=[E,E]\): the actual derived subgroup is perfect.
Its quotient by \(Z_E=E\cap Z\) is \(T^t\). Let \(H_i\) be the
inverse image in \(E\) of factor \(T_i\), and put \(E_i=H_i'\).
Again \(H_i=E_iZ_E\) and \(E_i\) is perfect. For \(i\ne j\),
\([H_i,H_j]\subset Z_E\). The commutator with a fixed element
of \(H_j\) is then a homomorphism \(H_i\to Z_E\), and kills
\(E_i\). Thus \(E_i\) commutes with \(H_j\), in particular with
\(E_j\). The product of the \(E_i\) has central quotient in
\(E\); perfection gives
\[
 E=E_1\cdots E_t,\qquad [E_i,E_j]=1\ (i\ne j),\qquad
 E_i/Z(E_i)\simeq T_i.
 \tag{G32.11}
\]
The assertion about the centre follows literally: a central
element of \(E_i\) also commutes with every other factor, hence
with \(H=EZ\), and is scalar by its absolute irreducibility.
Each \(E_i\) is therefore quasi-simple.

The representation of \(E\) is irreducible, since \(H=EZ\).
Every isotypic projector for \(E_i\) commutes both with \(E_i\)
and with all the other factors, hence with \(E\). It is therefore
scalar, so each restriction has only one isotypic type.
The commuting simple image algebras give the actual tensor product
\(V_{\mathbb C}=W_1\otimes\cdots\otimes W_t\), with each
\(\dim W_i>1\). One explicit map is obtained by choosing matrix
units \(e^{(i)}_{ab}\) in the actual commuting simple image
algebras. The products of their diagonal matrix units sum to
\(I\), and the joint off-diagonal matrix units identify all
their images. Thus the joint corner
\(\prod_i e^{(i)}_{11}V\) is nonzero. It is one line: a
non-scalar operator on a larger joint corner extends, through
the same joint matrix units, to a non-scalar commutant of
\(E\). Choose a nonzero vector \(u\) in that line.
Then
\[
 \bigotimes_i e_{a_i}\longmapsto
              \left(\prod_i e^{(i)}_{a_i1}\right)u
 \tag{G32.12}
\]
is the tensor isomorphism. Its inverse takes the coefficient
of \(u\) in \(\prod_i e^{(i)}_{1a_i}v\) for every multi-index;
the matrix-unit identities verify both composites. Nontriviality
of each factor and its perfection exclude a one-dimensional
factor representation.

Conjugation by \(N\) permutes the simple factors of \(H/Z\),
hence permutes the actual \(E_i\) and their matrix algebras.
The isomorphism (G32.12) is consequently a tensor decomposition
preserved by \(N\), including its factor permutations. HD-R14
excludes \(t>1\). We have obtained one actual quasi-simple group
\(S=E\), absolutely irreducible on the original space.

\paragraph{4. Characteristic core and the original normalizer.}
The minimal normal subgroup \(H/Z\) is unique. If there were a
distinct minimal normal subgroup with inverse image \(H_2\),
the two minimal normal subgroups would intersect trivially and
commute in \(N/Z\). Hence \([H_2,H]\subset Z\). As \(S=H'\)
is perfect, the same central-commutator homomorphism gives
\([H_2,S]=1\). The centralizer of the absolutely irreducible
\(S\) is scalar, so \(H_2\subset Z\), a contradiction.
Thus \(H\), and then its derived subgroup \(S\), are
characteristic in \(N\). Since \(N=N_3\) is normal in the original
\(G\), the original \(G\) normalizes \(S\).

Conjugation \(G\to\operatorname{Aut}(S)\) has kernel precisely
\(Q=G\cap K^*I\), by Schur. An automorphism of the perfect \(S\)
acting trivially on \(S/Z(S)\) has the form
\(s\mapsto s z_s\), where \(s\mapsto z_s\) is a homomorphism
into the centre. It is trivial by perfection. This proves the
injection (G32.1), and the image of the actual subgroup \(S\)
is the inner automorphism subgroup.

For each induced automorphism \(\alpha\), the actual operator
is retained through the finite \(K\)-linear equations
\[
                 T\rho(s)=\rho(\alpha(s))T\quad(s\in S).
 \tag{G32.13}
\]
It suffices to use any actual generating set of \(S\). Their
solution space has dimension one after extension to \(\mathbb C\)
by irreducibility, and contains the original normalizing matrix.
Thus it is a one-dimensional \(K\)-space, and every nonzero
solution is invertible. These are the actual outer-intertwiner
lines; no chosen extension character is substituted for them.

The faithful representation of \(S\) has all traces in \(K\).
Its centre is a cyclic scalar subgroup of \(\mu(K)\). Since
\(S\) is perfect, its determinant character is trivial, so a
central scalar \(z\) satisfies \(z^m=1\). This proves (G32.2).
In particular the centre has order at most two in the seven
fields with \(\mu(K)=\{\pm1\}\), and at most four in the Gaussian
field; its order also divides the displayed coefficient rank.

\paragraph{5. The actual coefficient factor and the remaining finite check.}
No curve or polarized factor has been replaced in this reduction:
\(S\), \(N\), their inclusions and all conjugating matrices above
are subgroups and operators on the original \(K^m\). The full
cover is still the original connected \(G\)-cover. Its prescribed
primitive matrix-coefficient idempotent remains the same element
of \(\mathbb Q[G]\), with its same integer-cleared Jacobian image,
restricted theta polarization and specified field action.

For a later choice of coordinates, a nonzero solution of the
finite intertwiner equations over \(K\) gives a literal map
\(C:K^m\to K^m\), and every original group matrix is transported
as \(g\mapsto C\rho(g)C^{-1}\), with the line \(Kw\) sent to
\(KCw\). Use the transported coefficient form, or compare it
with a prescribed positive invariant form \(h_0\). Absolute
irreducibility gives
\(h_0(Cx,Cy)=a h(x,y)\), \(a\in\mathbb Q_{>0}\).
Then the normalized coefficient
\(h_0(Cw,C\rho(g)w)/h_0(Cw,Cw)\) is exactly
\(h(w,\rho(g)w)/h(w,w)\). Thus the rational idempotent and actual
Jacobian image are identical. The Schur factors satisfy
\(c_0=c/a\), so \(c_0a=c\) on the original paired factor.
A change of primitive line, when needed, is the already displayed
rational matrix unit of HD-G31 with
\(b^\dagger b=(q_v/q_w)e_w\); its integer-cleared isogeny has
multiplier \(M^2q_v/q_w\). These are actual maps, not a
complex-character-only identification.

Hiss--Malle, \emph{Low-dimensional representations of quasi-simple
groups}, LMS J. Comput. Math.4 (2001),22--63, Theorem1 on p.22,
now applies to the actual faithful representation of \(S\).
The complete corrected finite tables are those in their
\emph{Corrigenda}, LMS J. Comput. Math.5 (2002),95--126;
the original generic families and the precise Table2/Table3
conventions in the 2001 paper, Section6 p.30, must also be retained.
The new reduction justifies using those quasi-simple tables in
this primitive branch; it does not perform their remaining
outer-normalizer checks.

The remaining check is finite in degrees15--35: retain the actual
faithful rows whose character fields embed in the specified \(K\)
and whose centre orders obey (G32.2), realize their \(K\)-matrix
representations, and analyze the actual finite normalizer matrices
in (G32.13) with the rank-three spectra of HD-G31. The full
original \(G\)-cover and its corner must remain linked by the
maps just given. Character-field or table matches alone do not
complete that check. This precise finite task remains open here;
the elementary-abelian normal case has been excluded in full.
